\title{xLSTM: Extended Long Short-Term Memory}

\begin{abstract}
During the 1990s, the concepts of the constant error carousel and gating emerged as fundamental components of the Long Short-Term Memory (LSTM) architecture. LSTMs have since proven robust and influential across a range of deep learning applications, notably serving as the backbone of the earliest Large Language Models (LLMs). Nonetheless, the introduction of Transformer models, characterized by inherently parallelizable self-attention mechanisms, has ushered in a paradigm shift, substantially surpassing LSTM performance at scale. This paper seeks to investigate an essential question: What are the limits of LSTM-based language modeling when scaled up to billions of parameters, enhanced by the latest advancements from modern LLM research, yet engineered to address LSTM's known constraints? To this end, we propose the use of exponential gating equipped with tailored normalization and stabilization methods. Furthermore, we reconfigure the LSTM's memory structure to create: (i) sLSTM, featuring scalar-valued memory, scalar-based updates, and an innovative memory mixing scheme, and (ii) mLSTM, a fully parallelizable variant with matrix-form memory and a covariance-driven update process. By embedding these LSTM refinements into residual block frameworks, we construct xLSTM modules which are then arranged in a residual stack to form comprehensive xLSTM models. The integration of exponential gating and these re-envisioned memory designs enables xLSTM to match or exceed the effectiveness and scalability of state-of-the-art Transformer and State Space Model counterparts.\\
Code available at: \url{https://github.com/NX-AI/xlstm}
\end{abstract}

\section{Introduction}
\label{sec:intro}

The Long Short-Term Memory (LSTM) ideas~\citep{Hochreiter:91,Hochreiter:97nips,Hochreiter:97}, i.e., the constant error carousel and gating, were introduced to overcome the vanishing gradient problem of recurrent neural networks~\citep{Hochreiter:91,Hochreiter:00book}:

\begin{align}
\label{eq:lstm_idea}
  c_t \ = \ f_t \ c_{t-1} \ + \ 
          i_t \  z_t \ , \quad  
          h_t \ = \ o_t \ \psi({c_t}) \ . 
\end{align}

The constant error carousel describes how the previous cell state $c_{t-1}$ (indicated in green) is incrementally adjusted with cell inputs $z_t$, a process regulated through the action of sigmoid gates (shown in blue). The input gate $i_t$ and forget gate $f_t$ are responsible for governing the cell state's update, whereas the output gate $o_t$ determines the final memory cell output, which is the hidden state $h_t$. Before being gated for output, the cell state undergoes normalization or squashing by $\psi$, resulting in the hidden state.

LSTMs have achieved notable success across a wide range of fields \citep{Hochreiter:01,Hochreiter:07,Schmidhuber:15}, dominating text generation tasks prior to the introduction of Transformers in 2017~\citep{Vaswani:17}. Their utility has been well established in numerous tasks involving sequences, including text synthesis~\citep{Graves:13,Karpathy:15blog}, handwriting generation~\citep{Graves:13}, sequence-to-sequence translation~\citep{Sutskever:14nips}, code evaluation~\citep{Zaremba:14arxiv}, image captioning~\citep{Karpathy:15,Hossain:19}, code generation~\citep{Karpathy:15blog}, as well as environmental modeling challenges such as rainfall-runoff simulations~\citep{Kratzert:18,Kratzert:19} and flood-warning hydrological modeling~\citep{Nearing:24}. In the context of reinforcement learning, LSTMs remain the superior sequence modeling choice, as demonstrated by applications like the AlphaStar model in StarCraft II~\citep{Vinyals:17short}, the OpenAI Five model for Dota 2~\citep{OpenAI:19}, and nuclear fusion magnetic control models~\citep{Degrave:22short}. LSTMs are particularly effective at learning abstract representations, skillfully capturing and retaining semantic content within their memory architecture~\citep{Karpathy:15blog}; evidence for this includes the discovery of units specializing in numbers and syntax~\citep{Lakretz:19}, language~\citep{Bau:19}, and sentiment~\citep{Radford:17arxiv}. LSTMs continue to be deployed in critical, real-world applications~\citep{Degrave:22short,Nearing:24}, demonstrating enduring relevance and resilience.

Although LSTMs have achieved remarkable success, they exhibit three central shortcomings: (i) Inability to update previously stored information. This issue is illustrated through the {\it Nearest Neighbor Search} problem (refer also to Appendix~\ref{sec:appNearestNeighborSearch}): Given a reference vector, the model must traverse a sequence in order to identify and return the value paired with the vector most similar to the reference at the sequence’s end. The leftmost graph in Figure~\ref{fig:lstmProblems} presents the mean squared error for this task. Standard LSTMs have difficulty replacing the current stored value when a more closely matching vector appears later, whereas our proposed xLSTM addresses this with exponential gating. (ii) Restricted memory capacity, necessitating compression of all retained information into scalar-valued cell states. We highlight this with {\it Rare Token Prediction}: as shown on the right in Figure~\ref{fig:lstmProblems}, LSTMs yield higher perplexity in predicting infrequent tokens from Wikitext-103~\citep{Merity:17} due to constrained memory resources. Our xLSTM, utilizing a matrix-based memory, mitigates this issue. (iii) Poor parallelizability caused by recurrent mixing of memory—hidden-to-hidden dependencies across time steps restrict LSTM processing to strictly sequential order.

These drawbacks have contributed to the rise of Transformers~\citep{Vaswani:17} for language modeling applications. An open question, then, is: What results could be achieved in language modeling by eliminating these limitations and scaling LSTM architectures up to the dimensions characteristic of modern Large Language Models?

\section{Extended Long Short-Term Memory}
\label{sec:xLSTM}

In order to address the shortcomings of LSTM, the Extended Long Short-Term Memory (xLSTM) framework makes two pivotal changes to the concept outlined in Equation~\eqref{eq:lstm_idea}. The innovations—namely exponential gating and distinctive memory architectures—lead to the creation of two new LSTM variants: (i) sLSTM (refer to Section~\ref{sec:sLSTM}), which employs a scalar memory, scalar updating, and incorporates memory mixing, and (ii) mLSTM (refer to Section~\ref{sec:mLSTM}), which utilizes a matrix-based memory system alongside a covariance (outer product) updating process that supports complete parallelization. Exponential gating serves as a common enhancement for both sLSTM and mLSTM. To facilitate parallel execution, mLSTM omits memory mixing, eliminating hidden-hidden recurrent connectivity. Both mLSTM and sLSTM models are scalable to multiple memory cells; in sLSTM, these cells interact through memory mixing. Moreover, sLSTM can operate with multiple heads—where mixing happens solely among the cells within a head, with no mixing between heads—thereby offering a new memory mixing paradigm when combined with exponential gating. For mLSTM, configurations with multiple heads and multiple cells are fundamentally interchangeable.

When these variants are placed within residual block modules, they form the xLSTM blocks, as detailed in Section~\ref{sec:xLSTMblock}. Layering these blocks residually in network designs yields xLSTM architectures, discussed in Section~\ref{sec:xLSTMarchitecure}. Figure~\ref{fig:xlstm_sketch} visually presents the components and structure of an xLSTM-based architecture.

\subsection{Review of the Long Short-Term Memory}
\label{sec:orgLSTM}

The original conception of LSTM~\citep{Hochreiter:91,Hochreiter:97nips,Hochreiter:97} established the scalar memory cell as a fundamental element for computation and storage, designed specifically to mitigate vanishing gradient issues~\citep{Hochreiter:91,Hochreiter:00book} via the constant error carousel mechanism—that is, through the way the cell state is updated. The memory cell architecture includes three distinct gates: the input gate, output gate, and forget gate, with the forget gate later added by \citet{Gers:00}. The temporal update equations governing the LSTM memory cell for time step $t$ are:

\begin{align}
c_t \ &= \  f_t \ c_{t-1} \ + \ i_t \ z_t &  & &\text{cell state} \\
h_t  \ &= \ o_t \ \tilde{h}_t \ , & \tilde{h}_t \ &= \ \psi \left( c_t\right)
  &\text{hidden state} \\
z_t \ &= \ \varphi \left( \tilde{z}_t \right) \ , 
  &\tilde{z}_t \ &=  \ \mathbf{w}^\top_{z} \ \mathbf{x}_t \ + \
  r_{z}  h_{t-1} \ + \  b_{z} \ \
  &\text{cell input} \\
i_t \ &= \ \sigma \left( \tilde{i}_t  \right) \ , 
  &\tilde{i}_t \ &= \ \mathbf{w}^\top_{i} \ \mathbf{x}_t \ + \
  r_{i}  \ h_{t-1} \ + \  b_{i} \ \
  &\text{input gate} \\
f_t \ &= \ \sigma \left(  \tilde{f}_t \right) \ , 
  &\tilde{f}_t \ &= \ \mathbf{w}^\top_{f} \ \mathbf{x}_t  \ + \
  r_{f}  \ h_{t-1} \ + \  b_{f} \ \
  &\text{forget gate} \\
o_t \ &= \ \sigma \left( \tilde{o}_t \right) \ , 
  &\tilde{o}_t  \ &= \ \mathbf{w}^\top_{o} \ \mathbf{x}_t \ + \
  r_{o}  \ h_{t-1} \ + \  b_{o} \ \
  &\text{output gate} 
\end{align}

The vectors $\mathbf{w}_{z}$, $\mathbf{w}_{i}$, $\mathbf{w}_{f}$, and $\mathbf{w}_{o}$ represent the input weight parameters that connect the input vector $\mathbf{x}_t$ to the cell input, the input gate, the forget gate, and the output gate, respectively. Likewise, $r_{z}$, $r_{i}$, $r_{f}$, and $r_{o}$ denote the recurrent weight parameters linking the previous hidden state $h_{t-1}$ to the cell input, input gate, forget gate, and output gate, in turn. The bias values associated with each component are $b_{z}$, $b_{i}$, $b_{f}$, and $b_{o}$. The activation functions for the cell input and the hidden state, $\varphi$ and $\psi$, respectively, are commonly chosen as $\tanh$. The hidden state activation $\psi$ ensures the cell state remains within a bounded range. The gate mechanisms employ sigmoid activations of the form $\sigma \left( x \right)= 1/(1+\exp(-x))$. Subsequent extensions to the architecture introduced memory cell vectors $\mathbf{c}_t \in \mathbb{R}^{d}$ composed of multiple scalar components, which facilitates the application of recurrent weight matrices $\mathbf{R} \in \mathbb{R}^{d \times d}$ to intermix the outputs from various memory cells~\citep{Greff:15}; refer to Appendix~\ref{sec:apporgLSTM} for further exposition. Empirical analyses through ablation studies have established that each memory cell element plays an essential role~\citep{Greff:15}.

\subsection{sLSTM}
\label{sec:sLSTM}

To enhance LSTM architectures with dynamic memory adjustment, we propose incorporating exponential gates (shown in red) along with mechanisms for normalization and stabilization. Specifically, exponential functions are utilized as activation functions for both input and forget gates. Regarding normalization, we define a normalizer state which accumulates the sum of the products of the input gate and all subsequent forget gates.

The scalar sLSTM forward pass is:

\begin{align}
\label{eq:slstmforward}
c_t \ &= \  f_t \ c_{t-1} \ + \ i_t \ z_t & & 
 &\text{cell state} \\
n_t \ &= \  f_t \ n_{t-1} \ + \ i_t  & & 
 &\text{normalizer state} \\
h_t \ &= \ o_t \ \tilde{h}_t \ ,
  &\tilde{h}_t \ &= \ c_t / n_t
&\text{hidden state} \\
z_t \ &= \ \varphi \left( \tilde{z}_t \right) \ , 
  &\tilde{z}_t \ &=  \ \mathbf{w}^\top_{z} \ \mathbf{x}_t \ + \
  r_{z}  h_{t-1} \ + \  b_{z}
  &\text{cell input} \\
i_t \ &= \ \!\exp \left( \tilde{i}_t  \right) \ , 
  &\tilde{i}_t \ &= \ \mathbf{w}^\top_{i} \ \mathbf{x}_t \ + \
  r_{i}  \ h_{t-1} \ + \  b_{i}
  &\text{input gate} \\
f_t \ &= \ \sigma \left(  \tilde{f}_t \right) \ \text{OR} \
\!\exp \left(  \tilde{f}_t \right) \ ,
  &\tilde{f}_t \ &= \ \mathbf{w}^\top_{f} \ \mathbf{x}_t  \ + \
  r_{f}  \ h_{t-1} \ + \  b_{f}
  &\text{forget gate} \\
o_t \ &= \ \sigma \left( \tilde{o}_t \right) \ , 
  &\tilde{o}_t  \ &= \ \mathbf{w}^\top_{o} \ \mathbf{x}_t \ + \
  r_{o}  \ h_{t-1} \ + \  b_{o}
  &\text{output gate} 
\end{align}

We transfer the original LSTM gating techniques, i.e., input- and/or hidden-dependent gating plus bias term, to the new architectures. Exponential activation functions can lead to large values that cause overflows. Therefore, we stabilize gates with an additional state \mbox{$m_t$~\citep{Milakov:18arxiv}}:

\begin{align}
\label{eq:slstmstabil}
m_t \ &= \ \max \left( \log ( f_t ) + m_{t-1} , \log ( i_t ) \right) &\text{stabilizer state} \\
i'_t \ &= \ \!\exp \left( \log \left ( i_t \right) - m_t \right) \ = \ \!\exp \left( \tilde{i}_t - m_t \right) \ \, 
  &\text{stabil. input gate} \\
f'_t \ &= \ \!\exp \left( \log \left( f_t \right) + m_{t-1} - m_t \right) \ \, 
  &\text{stabil. forget gate}
\end{align}

As demonstrated in Appendix~\ref{sec:appsLSTM}, substituting $f_t$ with $f'_t$ and $i_t$ with $i'_t$ during the forward pass leaves both the network's output and the gradients of the loss function with respect to its parameters unchanged.

\paragraph{Revised Memory Mixing.} Like its predecessor, the sLSTM architecture is capable of maintaining several memory cells (refer to Appendix~\ref{sec:appsLSTM}). Having multiple memory cells facilitates memory mixing via the recurrent connections $\mathbf{R}_{\mathbf{z}}$, $\mathbf{R}_{\mathbf{i}}$, $\mathbf{R}_{\mathbf{f}}$, and $\mathbf{R}_{\mathbf{o}}$, linking the hidden state $\mathbf{h}$ to the cell input $\mathbf{z}$ and the respective gates $\mathbf{i}$, $\mathbf{f}$, and $\mathbf{o}$. Exponential gating introduces a novel feature into memory mixing. The updated sLSTM design allows for multiple heads, enabling memory mixing inside individual heads, but not between heads. The pairing of exponential gating and the multi-head structure in sLSTM offers a distinct mechanism for memory mixing.

\subsection{mLSTM}
\label{sec:mLSTM}

To increase the capacity for information retention within LSTMs, we upgrade the traditional memory cell—originally a scalar $c \in \mathbb{R}$—to a matrix form $\mathbf{C} \in \mathbb{R}^{d \times d}$. This modification enables data retrieval through matrix multiplication. At a given time step $t$, we store a vector pair consisting of a key $\mathbf{k}_t \in \mathbb{R}^d$ and a value $\mathbf{v}_t \in \mathbb{R}^d$, following Transformer-based nomenclature. Subsequently, at a future time $t + \tau$, the value $\mathbf{v}_t$ is intended to be recovered using a query vector $\mathbf{q}_{t+\tau} \in \mathbb{R}^d$. This framework is consistent with the paradigm of Bidirectional Associative Memories (BAMs) \citep{Kohonen:72,Anderson:72,Nakano:72,Anderson:77}.

The covariance update rule~\citep{Sejnowski:77,Dayan:91} for storing a key-value pair is

\begin{align}
\mathbf{C}_t \ &= \ \mathbf{C}_{t-1} \ + \ \mathbf{v}_{t} \ \mathbf{k}_t^\top \ .
\end{align}

We apply layer normalization to the inputs prior to their projection as keys and values, ensuring that these projections have zero mean. According to~\citep{Dayan:91}, the covariance update rule yields optimality in terms of maximizing the separability among retrieved binary vectors, which directly relates to achieving the highest possible signal-to-noise ratio. Restricting retrieval operations to pairwise interactions and allowing quadratic computational complexity, as in attention mechanisms~\citep{Krotov:16,Krotov:17,Ramsauer:21}, can further increase separability. Furthermore, the covariance update procedure aligns with the approach used in Fast Weight Programmers~\citep{Schmidhuber:92ncfastweights,Schlag:21}, which were subsequently enhanced by introducing a fixed decay rate for $\mathbf{C}_{t-1}$ and a fixed learning rate for $\mathbf{v}_{t} \mathbf{k}_t ^\top$~\citep{Ba:16}. Following this rationale, we embed the covariance update scheme within the LSTM architecture, mapping the forget gate to the decay rate, the input gate to the learning rate, and assigning the scaling of the retrieved vector to the output gate.

For this matrix memory, the normalizer state is the weighted sum of key vectors, where each key vector is weighted by the input gate and all future forget gates. Again, the normalizer state keeps record of the strength of the gates. Since the dot product between query and normalizer state can be close to zero, we use the absolute value of this dot product and lower bound it by a threshold (typically 1.0) as done previously~\citep{Sun:23arxiv}. The mLSTM forward pass is:

\begin{align}
\label{eq:mlstm_recurrent_begin}
\mathbf{C}_t \ &= \  f_t \ \mathbf{C}_{t-1} \ + \ 
  i_t \ \mathbf{v}_t \ \mathbf{k}_t^\top & &  &\text{cell state} \\
\mathbf{n}_t \ &= \  f_t \ \mathbf{n}_{t-1} \ + \ 
  i_t \ \mathbf{k}_t & &  &\text{normalizer state} \\
\mathbf{h}_t  \ &= \ \mathbf{o}_t \ \odot \ \tilde{\mathbf{h}}_t
\ , \qquad \qquad 
\tilde{\mathbf{h}}_t \ = \   \mathbf{C}_t \mathbf{q}_t \ / \ 
  \max \left\{ |\mathbf{n}_t^\top \mathbf{q}_t|, 1 \right\} 
   &&&\text{hidden state} \\
\mathbf{q}_t \ &= \ \mathbf{W}_q \ \mathbf{x}_t \ + \ \mathbf{b}_q  & &  &\text{query input} \\
\mathbf{k}_t \ &= \ \frac{1}{\sqrt{d}} \mathbf{W}_k \ \mathbf{x}_t \ + \ \mathbf{b}_k  & &  &\text{key input} \\
\mathbf{v}_t \ &= \ \mathbf{W}_v \ \mathbf{x}_t \ + \ \mathbf{b}_v  & &  &\text{value input} \\
i_t \ &= \ \!\exp \!  \! \left( \tilde{i}_t  \right) \ , 
 \qquad \qquad \ \ \ \ \,
  \tilde{i}_t \ = \ \mathbf{w}^\top_{i} \ \mathbf{x}_t \ + \  b_{i} \ \ \ 
  &&&\text{input gate} \\
f_t \ &= \ \sigma \!  \left(  \tilde{f}_t \right) \ \text{OR} \ \!\exp \!  \! \left(  \tilde{f}_t \right) , \ \ \: \!
\tilde{f}_t \ = \ \mathbf{w}^\top_{f} \ \mathbf{x}_t  \ + \
  b_{f}
  &&& \text{forget gate} \\
\label{eq:mlstm_recurrent_end}
\mathbf{o}_t \ &= \ \sigma \left( \tilde{\mathbf{o}}_t \right) \ , \qquad \qquad \quad \ \ \ \,
  \tilde{\mathbf{o}}_t  \ = \ \mathbf{W}_{\mathbf{o}} \ \mathbf{x}_t \ + \
  \mathbf{b}_{\mathbf{o}}
  &&&\text{output gate} 
\end{align}

mLSTM is capable of supporting several memory cells, similar to the standard LSTM architecture. Within mLSTM, using multiple heads or multiple cells yields the same effect due to the absence of memory intermingling. To address the instability caused by mLSTM’s exponential gates, we employ stabilization strategies analogous to those applied in sLSTM, as detailed in Equation~\ref{eq:slstmstabil}. Because mLSTM does not perform memory mixing, its recurrence relationship lends itself to a parallelized formulation. Additional information can be found in Appendix~\ref{sec:appmLSTM}.

\subsection{xLSTM Architecture}
\label{sec:xLSTMarch}

\paragraph{xLSTM Blocks.}
\label{sec:xLSTMblock}
An xLSTM block aims to provide a nonlinear representation of previous inputs within a high-dimensional feature space. This transformation enhances the capacity to distinguish between distinct histories or contexts, which is critical for accurately forecasting the subsequent element in a sequence, such as the next token. Our approach draws upon Cover’s Theorem~\citep{Cover:65}, which posits that non-linear projections into higher-dimensional spaces increase the chances that patterns become linearly separable, compared to their separability in the original, lower-dimensional domain. We explore two residual block configurations: (i) A residual block with post up-projection (as in Transformers), wherein the historical sequence is first summarized non-linearly within the original space; afterward, a linear transformation projects this summary to a space of higher dimensionality, a non-linear activation is applied, and a final linear transformation returns it to the initial space. This mechanism is illustrated in the left side of Figure~\ref{fig:Backbones} and in the third column of Figure~\ref{fig:xlstm_sketch}. A comprehensive depiction can be found in Figure~\ref{fig:appsLSTM_detailed} in the appendix. (ii) A residual block utilizing pre up-projection (akin to State Space Models), in which the input first undergoes a linear mapping to a higher-dimensional space,

The past is aggregated in a high-dimensional manner through nonlinear operations, followed by a linear transformation that returns the representation to its original space. When an xLSTM block incorporates an sLSTM, we employ the up-projection block after the nonlinear operation. Conversely, in an xLSTM block with an mLSTM, the up-projection block is applied beforehand to accommodate increased memory capacity in the higher-dimensional setting. Further elaborations can be found in the left portion of Figure~\ref{fig:Backbones}, the third column of Figure~\ref{fig:xlstm_sketch}, or Figure~\ref{fig:appsLSTM_detailed} in the appendix.

\paragraph{xLSTM Architecture. }
\label{sec:xLSTMarchitecure}
To construct an xLSTM architecture, one stacks building blocks in a residual manner~\citep{Srivastava:15,He:16}. The backbone structure utilized follows the pre-LayerNorm approach~\citep{Ba:16b}, which is prevalent in state-of-the-art Large Language Models. Reference can be made to the last column of Figure~\ref{fig:xlstm_sketch} for visual detail.

\subsection{Memory and Speed Considerations}

In contrast to Transformers, xLSTM models exhibit linear computational complexity and maintain constant memory usage as the sequence length increases. The compressive memory architecture of xLSTM makes it particularly advantageous for use in industrial scenarios and on edge devices.

While mLSTM avoids parameterization of its memory, it incurs significant computational cost due to the use of a $d \times d$ memory matrix and a corresponding $d \times d$ update operation. This approach essentially balances larger memory capacity with increased computational effort. Despite this tradeoff, the required operations lend themselves well to parallelization on GPUs, resulting in minimal additional wall clock time.

Although mLSTM supports parallelism in a way that is comparable to FlashAttention~\citep{Dao:22,Dao:23} and GLA~\citep{Yang:23arxiv}, the sLSTM architecture cannot be parallelized efficiently as a consequence of its memory mixing via hidden-to-hidden interactions. To address this, we introduced an optimized CUDA implementation with memory handling down to the register level, achieving runtimes that are generally within a factor of two of mLSTM’s speed.

\section{Related Work}
\label{sec:related}

\paragraph{Linear Attention.} Various approaches have been proposed to address the quadratic time complexity of the Transformer with respect to context length, aiming to render attention mechanisms linear in this aspect. The Synthesizer generates synthetic attention weights independently of explicit token--token relationships~\citep{Tay:20arxiv}. Linformer achieves self-attention by employing a low-rank projection and enables linear approximation of the original operation~\citep{Wang:20arxiv}. Linear Transformer rewrites the standard attention mechanism in a form that allows for linear computation \citep{Katharopoulos:20}. Performer introduces a positive orthogonal random features method to provide a linearized softmax attention approximation~\citep{Choromanski:21}. Furthermore, alternative strategies replace the attention component with efficient long-sequence convolutions, as demonstrated in Structured Global Convolution (SGConv)~\citep{Li:22} and the Hyena Hierarchy~\citep{Poli:23}.

\paragraph{State Space Models.} State Space Models (SSMs) have gained significant attention recently due to their linear complexity with respect to sequence length and competitive results relative to Transformers. Early developments in this area included the Structured State Space sequence model (S4)~\citep{Gu:21}, which paved the way for successive approaches like the Diagonal State Space (DSS) model~\citep{Gupta:22}, Gated State Space (GSS) models~\citep{Mehta:22}, the S5 model~\citep{Smith:22}, Bidirectional Gated SSM (BiGS)~\citep{Wang:22}, H3 model~\citep{Fu:23}, as well as Mamba~\citep{Gu:24arxiv}.

\paragraph{Recurrent Neural Networks.} The linear scaling of context length has prompted proposals to substitute Transformers and attention mechanisms with Recurrent Neural Networks (RNNs). Deep Linear Recurrent Units (LRUs) in RNN architectures have demonstrated encouraging outcomes in language modeling tasks~\citep{Orvieto:23, De:24arxiv}. Similarly, variants such as the Hierarchically Gated Linear RNN (HGRN)~\citep{Qin:23} and its subsequent iteration HGRN2~\citep{Qin:24arxiv} have exhibited notable performance. Among approaches aimed at large-scale language modeling, RWKV~\citep{Peng:23arxivshort,Peng:24arxivshort} stands out, delivering results that are comparable with Transformer-based models.

\paragraph{Gating.}
Gating stands out as a fundamental concept introduced in LSTM, later adopted and reinterpreted across numerous contemporary methods. Approaches incorporating gating include HGRN~\citep{Qin:23}, HGRN2~\citep{Qin:24arxiv}, Gated Linear Attention (GLA)~\citep{Yang:23arxiv}, Gated State Space (GSS) models~\citep{Mehta:22}, Bidirectional Gated SSM (BiGS)~\citep{Wang:22}, Moving Average Equipped Gated Attention (MEGA)~\citep{Ma:22}, RWKV~\citep{Peng:23arxivshort}, as well as Mamba~\citep{Gu:24arxiv}.

\paragraph{Covariance Update Rule.} To increase memory storage efficiency, the mLSTM cell was augmented with a matrix-based memory following a covariance update rule. This technique echoes update mechanisms seen in Fast Weight Programmers \citep{Schmidhuber:92ncfastweights,Schlag:21}, RWKV-5 and RWKV-6~\citep{Peng:24arxivshort}, Retention~\citep{Sun:23arxiv}, Linear Transformer~\citep{Katharopoulos:20}, and HGRN2~\citep{Qin:24arxiv}.

\paragraph{Most Related.} The models that are most similar in concept to xLSTM include Retention~\citep{Sun:23arxiv}, RWKV~\citep{Peng:23arxivshort,Peng:24arxivshort}, and HGRN2~\citep{Qin:24arxiv}. These architectures incorporate principles such as matrix-based memory mechanisms and gating functions. Nevertheless, unlike the novel sLSTM, these methods lack support for memory mixing. The ability to mix memories is essential for addressing state tracking challenges, which makes LSTMs inherently more capable than both State Space Models (SSMs) and Transformers~\citep{Merrill:24,Deletang:23}. Effective state tracking is necessary for tasks such as code evaluation or maintaining information about entities across extended narratives.

\paragraph{Residually Stacking Architectures.} In line with virtually all recent large-scale deep neural networks, xLSTM architectures are developed through the residual stacking of modular units~\citep{Srivastava:15,He:16}.

This paradigm facilitated the evolution of very deep convolutional networks~\citep{He:16} as well as Transformers~\citep{Vaswani:17}. Transformers, in turn, serve as the foundation of Large Language Models (LLMs) such as GPT-3~\citep{Brown:20short}, ChatGPT~\citep{Schulman:22short}, GPT-4~\citep{Achiam:23gpt4short}, Megatron-LM~\citep{Shoeybi:19}, Gopher~\citep{Rae:21short}, ERNIE 3.0 Titan~\citep{Wang:21arxivshort}, GLaM ~\citep{Du:21glamshort}, Chinese M6~\citep{Lin:21}, mutilingual AlexaTM 20B~\citep{Soltan:22}, OPT~\citep{Zhang:22opt}, Chinchilla~\citep{Hoffmann:22short}, BLOOM~\citep{Scao:22arxivshort}, GLM-130B~\citep{Zeng:22arxivshort}, LaMDA~\citep{Thoppilan:22short}, PaLM~\citep{Chowdhery:22short}, Llama~\citep{Touvron:23arxiv}, and Gemini~\citep{GeminiTeam:23arxiv,Reid:24arxiv}.

\section{Experiments}
\label{sec:Exp}

This section presents a systematic empirical analysis of xLSTM in the context of language modeling, benchmarking it against established approaches. Section~\ref{sec:ExpSyntethic} explores xLSTM’s performance on artificial tasks designed to reveal its distinctive properties. In Section~\ref{sec:ExpComparison}, we report a comparison of validation set perplexity scores across several prominent language models, all trained on 15B tokens sampled from the SlimPajama dataset~\citep{Soboleva:23}. Ablation analyses for xLSTM are also conducted within this context. Subsequently, the scalability of each model is scrutinized following methodologies outlined in \citet{Kaplan:20} and \citet{Brown:20short}. 

In Section~\ref{sec:ExpLanguage}, we expand the evaluation, applying xLSTM and the most competitive models identified in Section~\ref{sec:ExpComparison} to a larger training corpus of 300B tokens from SlimPajama~\citep{Soboleva:23}. The extended study encompasses: an examination of generalization to longer input sequences, comparative validation perplexity and downstream task assessments as described by~\citep{Sutawika:24}, evaluation on 571 text domains from the PALOMA language benchmark~\citep{Magnusson:23arxivshort}, and finally, an investigation into scaling laws using a dataset twenty times the original size.

\label{sec:xlstm_blocks_notation}
In all our experiments, we denote the proportion of mLSTM to sLSTM blocks within xLSTM by xLSTM[$a$:$b$], signifying a ratio of $a/b$ for mLSTM-based to sLSTM-based blocks. For instance, xLSTM[7:1] indicates that, among eight total blocks, seven employ the mLSTM architecture while one utilizes sLSTM. In cases where the overall number of blocks is 48, this configuration results in 42 mLSTM blocks and 6 sLSTM blocks. Additionally, both pre up-projection and post up-projection blocks are incorporated for mLSTM and sLSTM, respectively, throughout all experimental conditions.

\subsection{Synthetic Tasks and Long Range Arena}
\label{sec:ExpSyntethic}
We begin by evaluating the impact of xLSTM's exponential gating mechanism incorporating memory mixing on formal language benchmarks~\citep{Deletang:23}. Next, we analyze how xLSTM's matrix memory influences results on the Multi-Query Associative Recall task~\citep{Arora:23arxiv}. Lastly, xLSTM's capabilities for handling extended sequences are measured using the Long Range Arena suite~\citep{Tay:21}.

\paragraph{Evaluation of xLSTM's Exponential Gating with Memory Mixing.} We evaluate the newly introduced exponential gating with memory mixing mechanism in xLSTM, designed to enhance its capacity for state tracking, as suggested in prior work~\citep{Merrill:24, Merrill:23}. For our study, we adopt and expand upon the suite of formal language tasks provided by \citet{Deletang:23}, enabling variable-length sequence training to support evaluation of length generalization. The complete set of task descriptions and additional results are provided in Appendix~\ref{sec:appExpSynthetic-formalLang}. 

Our evaluation contrasts xLSTM against a range of baseline methods, including variants of Transformers, State Space Models, and various forms of Recurrent Neural Networks. Performance is measured by accuracy on the subset of tokens essential for the tasks; accuracy values range from 0 (corresponding to random guessing) to 1 (denoting perfect prediction). We consider 2-block configurations for the competing methods: xLSTM[0:1] (corresponding to an sLSTM only), xLSTM[1:0] (mLSTM only), xLSTM[1:1], Llama, Mamba, RWKV, Retention, Hyena, LSTM, as well as LSTM implemented within Transformer architectures (LSTM (Block)). Figure~\ref{fig:formal-main} presents a summary of the empirical findings. 

We observe that models such as Transformers and State Space Models, when lacking memory mixing (and thus state tracking), are incapable of solving certain types of regular grammar problems—for instance, the parity task. This observation supports previous conclusions that, in terms of computational power, Transformers and State Space Models are fundamentally less expressive than RNNs~\citep{Merrill:24, Merrill:23, Deletang:23}.

\paragraph{Test of xLSTM's Memory Capacities on Associative Recall Tasks.} This study evaluates the matrix memory of xLSTM by measuring its memory capacity on the Multi-Query Associative Recall task~\citep{Arora:23arxiv}. In this task, sequences are constructed such that each one contains key-value pairs, randomly sampled from an extensive vocabulary, which the models are required to retain for subsequent retrieval. To further increase the challenge relative to previous setups, we scale up the number of key-value pairs to as many as 256 and lengthen the context up to 2048 tokens. This expanded benchmark enables a more thorough assessment of the memory capacities across various models. We compare 2-block variants of Llama, Mamba, RWKV-5, RWKV-6, xLSTM[1:1], and xLSTM[1:0], using recall accuracy as the primary evaluation metric. Due to the fact that Transformers like Llama possess a memory capacity that grows exponentially with respect to their coding dimension~\citep{Ramsauer:21}, they serve as the performance benchmark for this task. The outcomes are visualized in Figure~\ref{fig:mqar-main}. 
Among non-Transformer models, xLSTM[1:1] consistently achieves the highest recall accuracy, outperforming others even in smaller model configurations. Notably, the inclusion of the sLSTM block in xLSTM does not reduce memory capacity; on the contrary, it augments it—an effect most pronounced in the hardest setting with 256 key-value pairs. Further experimental details are included in Appendix~\ref{sec:appExpSynthetic-mqar}, where extrapolation studies show that xLSTM's superior memory properties persist in scenarios with longer contexts than those encountered in training.

\paragraph{Test of xLSTM's Long Context Capabilities on Long Range Arena.} In order to evaluate how xLSTM manages extended sequences and substantial contexts, we perform a comparative analysis of various approaches on the Long Range Arena~\citep{Tay:21}. Across all tasks, xLSTM reliably achieves high results, indicating that its architecture is notably adept at addressing diverse challenges associated with long-context scenarios.

For more details, see Appendix~\ref{sec:appExpSynthetic-lra}.

\subsection{Method Comparison and Ablation Study}
\label{sec:ExpComparison}

To investigate the central inquiry posed by this study—specifically, how well the proposed LSTM extensions perform when applied at scale in language modeling tasks—we conduct experiments involving xLSTMs, Transformers, State Space Models, and additional approaches. Each model is trained using 15B tokens extracted from SlimPajama under consistent auto-regressive conditions. Performance is assessed through comparisons on a shared validation set, and further analysis is carried out via ablation studies focusing on the xLSTMs.

\paragraph{Comparing xLSTM to Other Methods.} To facilitate a comparative analysis, we train each model using 15B tokens drawn from SlimPajama~\citep{Soboleva:23} and assess performance through perplexity on the corresponding validation set. The set of baselines includes xLSTM (the novel architecture proposed here), as well as GPT-3 (Transformer)~\citep{Brown:20short}, Llama (Transformer)~\citep{Touvron:23arxiv}, H3 (SSM)~\citep{Fu:23}, Mamba (SSM)~\citep{Gu:24arxiv}, RWKV-4, RWKV-5, RWKV-6 (each RNNs)~\citep{Peng:23arxivshort, Peng:24arxivshort}, GLA (linear Transformer)~\citep{Yang:23arxiv}, HGRN, HGRN2 (RNNs)~\citep{Qin:23, Qin:24arxiv}, RetNet (linear Transformer)~\citep{Sun:23arxiv}, Hyena (linear Transformer)~\citep{Poli:23}, as well as the xLSTM[1:0] and xLSTM[7:1] variants (refer to Section~\ref{sec:xlstm_blocks_notation}). With regard to training configuration, all models utilize mixed precision. However, for RWKV-5, RWKV-6, GLA, and HGRN2, the PyTorch automated mixed precision functionality was not employed (see additional details in Appendix Section~\ref{sec:appGenTrainProc}).

These architectures fall into three main categories: (a) Transformers, (b) State Space Models (SSMs), and (c) RNNs grouped alongside linear Transformers—where the latter refers to models that replace standard attention mechanisms with linear operations. Each network is constrained to be approximately the size of a GPT-3 model containing 350M parameters, defined by a 1024-dimensional embedding space and 24 residual layers. Notably, only GPT-3 shares parameters between input and output token embeddings, resulting in a slightly reduced parameter count. Empirical results presented in Table~\ref{tab:spaj15b_model_results} indicate that xLSTM consistently achieves superior perplexity on the validation set relative to all benchmarks. Additional details may be found in Appendix~\ref{sec:appExpComparison}. Furthermore, Figure~\ref{fig:temp_sclaw15B} illustrates the observed scaling trends, suggesting that xLSTM retains its competitive edge as model size increases.

\paragraph{Ablation Studies.}
In Table~\ref{tab:spaj15b_model_results} and Figure~\ref{fig:temp_sclaw15B}, xLSTM displays strong performance in language modeling tasks when trained with 15B tokens from SlimPajama. To systematically dissect the architectural modifications leading from the original LSTM to xLSTM, we incrementally transform a standard LSTM framework toward the xLSTM design. We begin by incorporating LSTM layers within pre-LayerNorm residual structures. The next phase involves extending the architecture to include a post up-projection block. The final modification introduces exponential gating together with matrix memory components.

The results are shown in Table~\ref{tab:ablstudies} (top). 

Ablation experiments suggest that the notable gains in performance stem from both the use of exponential gating and the incorporation of matrix memory. Furthermore, recognizing the critical role of gating within RNNs and State Space Models, we investigate the effects of varying gating strategies. As shown in Table~\ref{tab:ablstudies} (bottom), our findings indicate that allowing each gate to be trainable and responsive to the input further enhances performance. More detailed ablations focusing on the separate backbone elements can be found in Appendix~\ref{sec:appAblDetails}.

\subsection{xLSTM as Large Language Model}
\label{sec:ExpLanguage}

Our investigation concludes with extensive experiments in language modeling at scale, aiming to assess the suitability of xLSTM as a large language model (LLM). To this end, we expand the training dataset size and conduct training on 300B tokens sourced from SlimPajama. This volume of training data aligns with that used by models such as Mamba~\citep{Gu:24arxiv} and Griffin~\citep{De:24arxiv}.

For comparison, xLSTM is evaluated against RWKV-4, Llama, and Mamba, with each model representing the respective approach classes delineated in Section~\ref{sec:ExpComparison}.
RWKV-4 is chosen as the RNN benchmark since appropriate training precision configurations for RWKV-5, RWKV-6, and HGRN2 became available only after commencing the 300B token training phase (refer to Appendix~\ref{sec:appGenTrainProc}).
We train models with varied sizes (125M, 350M, 760M, and 1.3B), examine their capacity for length generalization, and validate their performance using the validation dataset. In addition, we evaluate these models on downstream task metrics, gauge their proficiency in language modeling across 571 distinct text domains covered by the PALOMA benchmark, and analyze their scaling law characteristics.

\paragraph{Sequence Length Extrapolation.}
\label{sec:spaj300B_ctx_extrapolation}
To begin, we assess the extrapolation capabilities across sequence lengths for large-scale (1.3B parameters) variants of xLSTM, RWKV-4, Llama, and Mamba. Each model is trained using sequences of length 2048 and subsequently evaluated on sequences extending up to length 16384. The resulting performance is visualized in Figure~\ref{fig:spaj300B_seq_extrapolation}. xLSTM distinctly stands out by retaining lower perplexity scores over these extended sequence lengths compared to its counterparts.

\paragraph{Validation Perplexity and Downstream Tasks.}
\label{sec:spaj_downstream_eval}
Next, we systematically measure the abilities of xLSTM, RWKV-4, Llama, and Mamba, spanning all available parameter scales, on the SlimPajama validation benchmark for next-token prediction. Furthermore, these models are also assessed on a series of downstream benchmarks, each designed to evaluate their common sense reasoning competencies.

The third column in Table~\ref{tab:lmeval} presents the perplexity scores on the validation set for each approach. Notably, xLSTM[1:0] and xLSTM[7:1] consistently achieve the lowest validation perplexity across all model scales, establishing them as the top performers. The remaining columns in Table~\ref{tab:lmeval} summarize results for a variety of downstream tasks. xLSTM generally outperforms other methods on nearly all tasks and sizes, with the exception of the ARC task, where Mamba occasionally demonstrates superior performance. Further details can be found in Appendix~\ref{sec:appExpLanguage}.

\paragraph{Performance on PALOMA Language Tasks.} Additionally, we evaluate the next token prediction capabilities of xLSTM, RWKV-4, Llama, and Mamba across all available model sizes using PALOMA language tasks~\citep{Magnusson:23arxivshort}. For this assessment, we utilize perplexity as the metric, calculated for each of the 571 distinct text domains, which span sources from nytimes.com to Reddit’s r/depression.

As summarized in Table~\ref{tab:ppleval}, the results are partitioned into language modeling (the initial seven columns) and more specific domain benchmarks (the remaining five columns). Among the xLSTM variants, xLSTM[1:0] demonstrates superior results over xLSTM[7:1] on these language tasks.

When compared across domains, xLSTM[1:0] yields lower perplexity than Mamba in 568 out of 571 (99.5\%) cases, outperforms Llama in 486 of 571 (85.1\%) domains, and achieves better perplexity than RWKV-4 in 570 of 571 (99.8\%) instances. Further elaboration is provided in Appendix~\ref{sec:appExpLanguage}.

\paragraph{Scaling Laws.} Next, we investigate the characteristic power-law scaling relationship, which facilitates predicting model performance as the parameter count increases~\citep{Kaplan:20,Brown:20short}. The scaling curves are illustrated in Figure~\ref{fig:temp_sclaw300B}. While all examined architectures exhibit comparable scaling trajectories, the results are shifted by distinct baseline values. RWKV-4 yields the lowest performance among the tested models, followed sequentially by Llama and Mamba. xLSTM achieves improved results relative to Mamba, displaying a gap similar to that between Mamba and Llama. These scaling trends suggest that with further model enlargement, xLSTM is likely to retain its advantage over Transformer and State-Space model variants.

\textbf{Generation Times and Maximal Throughput.} In Figure~\ref{fig:inference_time_generation}, we present measurements of text generation latency, while maximal throughput is illustrated in Figure~\ref{fig:inference_time_decoding_throughput} (left), evaluating xLSTM variants at the 1.3B parameter level. Our comparison group consists of equivalently-sized implementations of Mamba, Llama, and RWKV sourced from HuggingFace, where the Llama model is additionally benchmarked with a static key-value cache. During experimentation, the compilation of the Transformer Model's full cache and Mamba model via \texttt{torch.compile} were not functional. In these generation benchmarks, all models were evaluated with a batch size of 1 and a pre-fill of 16—settings that should optimally benefit the Transformer architecture.

The results in Figure~\ref{fig:inference_time_generation} highlight the linear growth in generation time for xLSTM, Mamba, and RWKV-4, as opposed to the quadratic complexity observed with Llama. When assessing decoding throughput, we further explore varying batch sizes and pre-fill configurations for the Llama baseline. As shown in Figure~\ref{fig:inference_time_decoding_throughput} (right), xLSTM exploits its fixed memory requirements to accommodate significantly larger batch sizes relative to Llama, thereby delivering the greatest throughput among the models evaluated.

\section{Limitations}

(i) Unlike mLSTM, the memory mixing mechanism in sLSTM necessitates sequential operations, preventing efficient parallelization and consequently hindering fast parallel execution. However, we have implemented a high-speed CUDA kernel tailored for sLSTM, and it currently operates at under twice the duration of our parallel mLSTM kernel. (ii) The existing CUDA kernels for mLSTM have not undergone optimization, resulting in a performance that is approximately four times slower compared to FlashAttention or the scan approach adopted by Mamba. It is plausible that substantially faster CUDA kernels can be realized, analogous to the strategies employed in FlashAttention. (iii) Processing matrix memory in mLSTM involves significant computational demands due to the requirement of handling $d \times d$ matrices. Nevertheless, since parameter-free memory update and retrieval can leverage parallelized matrix operations, the impact of this complex memory mechanism on overall wall clock time remains minimal. (iv) Careful selection of initial values for the forget gates is essential. (v) Although the matrix memory of mLSTM does not scale with sequence length, increasing the context size could place greater strain on memory capacity. To date, this has not been a constraint for contexts reaching up to 16k tokens, as detailed in Section~\ref{sec:spaj300B_ctx_extrapolation}. (vi) Due to the considerable computational requirements associated with large-scale language experiments, neither the architecture nor the hyperparameter settings have been comprehensively optimized, especially for the larger variants of xLSTM. We expect that a thorough and systematic optimization effort will be necessary for xLSTM to fully achieve its capabilities.

\section{Conclusion}
We have provided a partial answer to our primary research question: To what extent can LSTM scale up to billions of parameters in the context of language modeling? Based on our results, it is clear that LSTM architectures can achieve performance comparable to contemporary approaches such as Transformers and State Space Models. Through improvements—specifically, the introduction of exponential gating with memory mixing and a novel memory design—we have evolved LSTM into xLSTM. xLSTM demonstrates strong results in language modeling, outperforming or matching the latest Transformer and State Space Model methods. The observed scaling trends suggest that larger xLSTM variants could emerge as strong contenders against existing Transformer-based Large Language Models. Furthermore, xLSTM shows promise for broad applications beyond language modeling, including domains such as Reinforcement Learning, Time Series Forecasting, and modeling dynamical physical phenomena.

\end{document}